% ANCHOR: after_label def:c2-even
\begin{remark}[Physics notation: conformal and R-symmetry generators]
The even generators are the vector fields $v^m(x)\partial_m$, with $X_{mn}=\partial_mv_n-\partial_nv_m$:
\begin{itemize}
\item translation $P_b$: $v^m=b^m$, so $P_m=\partial_m$;
\item rotation with $\omega_{mn}=-\omega_{nm}$: $v^m=\omega_{mn}x_n$, $X_{mn}=-2\omega_{mn}$, the Euclidean Lorentz algebra $\mathfrak{so}(4)=\mathfrak{su}(2)_L\oplus\mathfrak{su}(2)_R$;
\item dilatation $\Delta=2E$: $v^m=2x^m$, $\operatorname{div}v=8$;
\item special conformal $K_b$: $v^m=x^2b^m-2(b\cdot x)x^m$, $X_{mn}=4(x_mb_n-x_nb_m)$, $\operatorname{div}v=-8\,b\cdot x$;
\item R-symmetry $R\in\mathfrak{so}(6)$: $v(y)^I=R_{IJ}y_J$ on $W_{\mathbb R}^\perp$, acting on $W$ by the traceless $r^i{}_j=\tfrac14R_{IJ}(\bar\Sigma^J\Sigma^I)^i{}_j$ with $\bar\Sigma^I=(\Sigma^I)^\dagger$, an element of the complexified $\mathfrak{su}(4)_R$.
\end{itemize}
The even part is $\mathfrak{sl}(4,\mathbb C)\oplus\mathfrak{sl}(4,\mathbb C)$, the complexification of the conformal algebra $\mathfrak{su}(2,2)$ and of $\mathfrak{su}(4)_R$.
\end{remark}

% ANCHOR: after_label def:c2-odd
\begin{remark}[Physics notation: conformal Killing spinors]
With $Q_1=(\xi_\alpha{}^i,\bar\xi_{\dot\alpha i})$, the parameters of $Q^i_\alpha$ and $\bar Q_{\dot\alpha i}$, and $Q_2=(\eta^\alpha{}_i,\bar\eta^{\dot\alpha i})$, the parameters of $S^\alpha{}_i$ and $\bar S^{\dot\alpha i}$, the conformal Killing spinor $x\mapsto Q_1+\rho(x)Q_2$ is
\[
\epsilon_\alpha{}^i(x)=\xi_\alpha{}^i-i\,x_m\sigma^m_{\alpha\dot\alpha}\bar\eta^{\dot\alpha i},\qquad \bar\epsilon_{\dot\alpha i}(x)=\bar\xi_{\dot\alpha i}+i\,x_m\eta^\alpha{}_i\sigma^m_{\alpha\dot\alpha}:
\]
the parameter $\bar\eta$ of $\bar S$ enters the undotted $\epsilon$, and the parameter $\eta$ of $S$ the dotted $\bar\epsilon$. The components $\xi$, $\bar\xi$, $\eta$, $\bar\eta$ are commuting numbers; the odd parity of the generators sits in $\Pi$.
\end{remark}

% ANCHOR: after_label def:c2-bracket
\begin{remark}[Physics notation: the brackets of the superconformal algebra]
For $\epsilon_a=\xi_a+\rho(x)\eta_a$, $a=1,2$,
\[
\begin{aligned}
&[\epsilon_1,\epsilon_2]=\Big(x\mapsto i\big(\bar\epsilon_{1\dot\alpha i}\bar\sigma^{m\dot\alpha\alpha}\epsilon_{2\alpha}{}^i+\bar\epsilon_{2\dot\alpha i}\bar\sigma^{m\dot\alpha\alpha}\epsilon_{1\alpha}{}^i\big)\partial_m\Big)+r,\\
&r^k{}_i=2\big(\xi_{1\alpha}{}^k\eta_2{}^\alpha{}_i+\bar\xi_{1\dot\alpha i}\bar\eta_2^{\dot\alpha k}\big)+2\big(\xi_{2\alpha}{}^k\eta_1{}^\alpha{}_i+\bar\xi_{2\dot\alpha i}\bar\eta_1^{\dot\alpha k}\big)-\mathrm{trace},
\end{aligned}
\]
with the rotation of $W_{\mathbb R}^\perp$ written as the $SU(4)_R$ rotation $r^k{}_i$ of $W$. For an even generator with $v^m$, $X_{mn}$, $\operatorname{div}v$ and $r^i{}_j$ as after Definition \ref{def:c2-even},
\[
\begin{aligned}
[v,\epsilon]_\alpha{}^i&=v^n\partial_n\epsilon_\alpha{}^i-\tfrac18X_{mn}(\sigma^m\bar\sigma^n)_\alpha{}^\beta\epsilon_\beta{}^i-\tfrac18\operatorname{div}v\,\epsilon_\alpha{}^i+r^i{}_j\epsilon_\alpha{}^j,\\
[v,\epsilon]_{\dot\alpha i}&=v^n\partial_n\bar\epsilon_{\dot\alpha i}+\tfrac18X_{mn}\bar\epsilon_{\dot\beta i}(\bar\sigma^m\sigma^n)^{\dot\beta}{}_{\dot\alpha}-\tfrac18\operatorname{div}v\,\bar\epsilon_{\dot\alpha i}-\bar\epsilon_{\dot\alpha j}r^j{}_i ,
\end{aligned}
\]
so $\tfrac14\rho(d(v,-))$ is the spin term $-\tfrac18X_{mn}\sigma^m\bar\sigma^n$ and $-\tfrac18\operatorname{div}v=-1$ for $\Delta=2E$. For a rotation the spin terms of $[v,\epsilon]_\alpha{}^i$ and $[v,\epsilon]_{\dot\alpha i}$ are
\[
\tfrac14\omega_{mn}(\sigma^m\bar\sigma^n)_\alpha{}^\beta\epsilon_\beta{}^i,\qquad -\tfrac14\omega_{mn}\bar\epsilon_{\dot\beta i}(\bar\sigma^m\sigma^n)^{\dot\beta}{}_{\dot\alpha}.
\]
\end{remark}

% ANCHOR: end_of_section 2.1
\begin{remark}[Physics notation: named generators and their brackets]
Under $SU(2)_L\times SU(2)_R\times SU(4)_R$ the named subspaces are: $P_m$ and $K_m$ in $(\mathbf 2,\mathbf 2,\mathbf 1)$, the rotations in $(\mathbf 3,\mathbf 1,\mathbf 1)\oplus(\mathbf 1,\mathbf 3,\mathbf 1)$, $\Delta$ a singlet, the R-symmetry in $(\mathbf 1,\mathbf 1,\mathbf{15})$, $Q^i_\alpha$ in $S_+\otimes W=(\mathbf 2,\mathbf 1,\mathbf 4)$, $\bar Q_{\dot\alpha i}$ in $S_-\otimes W^*=(\mathbf 1,\mathbf 2,\bar{\mathbf 4})$, $\bar S^{\dot\alpha i}$ in $S_-\otimes W=(\mathbf 1,\mathbf 2,\mathbf 4)$ and $S^\alpha{}_i$ in $S_+\otimes W^*=(\mathbf 2,\mathbf 1,\bar{\mathbf 4})$. In the normalisation $\Delta=-E$ of [SW, (3)] the eigenvalues of $\mathrm{ad}_\Delta$ are the scaling dimensions $1$ of $P$, $\tfrac12$ of $Q,\bar Q$, $-\tfrac12$ of $S,\bar S$ and $-1$ of $K$. The brackets of Proposition \ref{prop:c2-special} read:
\begin{itemize}
\item $Q_1,Q_2\in\Sigma$: the translation $v^m=i\big(\bar\xi_{2\dot\alpha i}\bar\sigma^{m\dot\alpha\alpha}\xi_{1\alpha}{}^i+\bar\xi_{1\dot\alpha i}\bar\sigma^{m\dot\alpha\alpha}\xi_{2\alpha}{}^i\big)$, pairing $Q^i_\alpha$ with $\bar Q_{\dot\alpha i}$;
\item $Q_1,Q_2\in\Sigma^*$: $-K_{v_1}$ with $v_1^m=-i\big(\eta_2{}^\alpha{}_i\sigma^m_{\alpha\dot\alpha}\bar\eta_1^{\dot\alpha i}+\eta_1{}^\alpha{}_i\sigma^m_{\alpha\dot\alpha}\bar\eta_2^{\dot\alpha i}\big)$, pairing $S^\alpha{}_i$ with $\bar S^{\dot\alpha i}$;
\item $Q_1\in\Sigma$, $Q_2\in\Sigma^*$: the dilatation part $\tfrac12\big(\xi_\alpha{}^i\eta^\alpha{}_i-\bar\xi_{\dot\alpha i}\bar\eta^{\dot\alpha i}\big)\Delta$, a rotation, and the R-symmetry rotation $r$;
\item translations: $[b,\eta]=\rho(b)\eta$ with $(\rho(b)\eta)_\alpha{}^i=-i\,b_m\sigma^m_{\alpha\dot\alpha}\bar\eta^{\dot\alpha i}$, from $\bar S$ to $Q$, and $(\rho(b)\eta)_{\dot\alpha i}=i\,b_m\eta^\alpha{}_i\sigma^m_{\alpha\dot\alpha}$, from $S$ to $\bar Q$;
\item special conformal: $[K_b,\xi]=\rho(b)\xi$ with $(\rho(b)\xi)^{\dot\alpha i}=-i\,b_m\bar\sigma^{m\dot\alpha\alpha}\xi_\alpha{}^i$, from $Q$ to $\bar S$, and $(\rho(b)\xi)^\alpha{}_i=i\,b_m\bar\xi_{\dot\alpha i}\bar\sigma^{m\dot\alpha\alpha}$, from $\bar Q$ to $S$;
\item $[\Delta,Q_1]=-Q_1$ on $Q^i_\alpha$, $\bar Q_{\dot\alpha i}$ and $[\Delta,Q_2]=+Q_2$ on $S^\alpha{}_i$, $\bar S^{\dot\alpha i}$, for $\Delta=2E$.
\end{itemize}
\end{remark}

% ANCHOR: end_of_section 2.2
\begin{remark}[Physics notation: constant ghosts]
A cochain $f(x_1,\dots,x_k)$ is the coefficient of $k$ constant ghosts of $\mathfrak{psl}(4|4)$, one for each argument, all of ghost number $+1$: the ghosts of $P_m$, $K_m$, $\Delta$, the rotations and the R-symmetry are Grassmann odd, and those of $Q^i_\alpha$, $\bar Q_{\dot\alpha i}$, $S^\alpha{}_i$, $\bar S^{\dot\alpha i}$ are the commuting spinors $\xi_\alpha{}^i$, $\bar\xi_{\dot\alpha i}$, $\eta^\alpha{}_i$, $\bar\eta^{\dot\alpha i}$. An element of degree $0$ has total ghost number $0$, so $S^{(k)}$ has ghost number $-k$, carried by its antifields. $d_{CE}$ is the BRST differential of the constant ghosts, the two sign conventions of Lemma \ref{lem:c2-dcesign} differ by the sign of every constant ghost, and \eqref{eq:c2-mc} is the classical master equation of $S+S^{(1)}+S^{(2)}$ extended by the constant ghosts.
\end{remark}

% ANCHOR: after_label def:c2-S1
\begin{remark}[Physics notation: conformal and R-symmetry transformations of the fields]
Every $S^{(1)}(x_1)$ has the form
\[
S^{(1)}(x_1)=\int d^4x\Big\{(x_1.A_m,A^*_m)+\tfrac14(x_1.A_{ij},A^{*ij})-(x_1.\lambda_\alpha{}^i,\lambda^{*\alpha}{}_i)+(x_1.\bar\lambda_{\dot\alpha i},\bar\lambda^{*\dot\alpha i})-(x_1.c,c^*)\Big\},
\]
with $x_1.\varphi$ the transformation of the field $\varphi$ along $x_1$. For an even generator with $v^m$, $X_{mn}$, $\operatorname{div}v$ and $r^i{}_j$ as after Definition \ref{def:c2-even},
\[
\begin{aligned}
v.A_m&=v^n\partial_nA_m+A_n\partial_mv^n,\qquad v.A_{ij}=v^n\partial_nA_{ij}+\tfrac14\operatorname{div}v\,A_{ij}-r^k{}_iA_{kj}-r^k{}_jA_{ik},\qquad v.c=v^n\partial_nc,\\
v.\lambda_\alpha{}^i&=v^n\partial_n\lambda_\alpha{}^i-\tfrac18X_{mn}(\sigma^m\bar\sigma^n)_\alpha{}^\beta\lambda_\beta{}^i+\tfrac38\operatorname{div}v\,\lambda_\alpha{}^i+r^i{}_j\lambda_\alpha{}^j,\\
v.\bar\lambda_{\dot\alpha i}&=v^n\partial_n\bar\lambda_{\dot\alpha i}+\tfrac18X_{mn}\bar\lambda_{\dot\beta i}(\bar\sigma^m\sigma^n)^{\dot\beta}{}_{\dot\alpha}+\tfrac38\operatorname{div}v\,\bar\lambda_{\dot\alpha i}-\bar\lambda_{\dot\alpha j}r^j{}_i :
\end{aligned}
\]
the conformal transformations of a vector, of six scalars of Weyl weight $1$ and of Weyl spinors of Weyl weight $\tfrac32$, the coefficients $\tfrac14$ and $\tfrac38$ of $\operatorname{div}v$ being the weight divided by $4$, together with the $SU(4)_R$ rotation $r$.
\end{remark}

% ANCHOR: end_of_section 2.3
\begin{remark}[Physics notation: supersymmetry and special superconformal transformations]
For $\xi\in\Sigma$, \eqref{eq:c2-S1Q} gives $\xi.A=-\Gamma(\xi,\lambda)$ and $\xi.\lambda=-\tfrac12\rho(F_A)\xi$, in components
\[
\begin{aligned}
\xi.A_m&=-i\big(\bar\lambda_{\dot\alpha i}\bar\sigma^{m\dot\alpha\alpha}\xi_\alpha{}^i+\bar\xi_{\dot\alpha i}\bar\sigma^{m\dot\alpha\alpha}\lambda_\alpha{}^i\big),\qquad
\xi.A_{kl}=-2\varepsilon_{jikl}\xi_\alpha{}^j\varepsilon^{\alpha\beta}\lambda_\beta{}^i+2\big(\bar\xi_{\dot\alpha k}\varepsilon^{\dot\alpha\dot\beta}\bar\lambda_{\dot\beta l}-\bar\xi_{\dot\alpha l}\varepsilon^{\dot\alpha\dot\beta}\bar\lambda_{\dot\beta k}\big),\\
\xi.\lambda_\alpha{}^i&=\tfrac14F_{mn}(\sigma^m\bar\sigma^n)_\alpha{}^\beta\xi_\beta{}^i+\tfrac14[A^{ik},A_{kj}]\,\xi_\alpha{}^j+\tfrac i2\,\sigma^m_{\alpha\dot\alpha}\varepsilon^{\dot\alpha\dot\beta}D_mA^{ij}\,\bar\xi_{\dot\beta j},\\
\xi.\bar\lambda_{\dot\alpha i}&=-\tfrac14F_{mn}\bar\xi_{\dot\beta i}(\bar\sigma^m\sigma^n)^{\dot\beta}{}_{\dot\alpha}+\tfrac14[A_{ik},A^{kj}]\,\bar\xi_{\dot\alpha j}-\tfrac i2\,\varepsilon^{\alpha\beta}\sigma^m_{\alpha\dot\alpha}D_mA_{ij}\,\xi_\beta{}^j ,
\end{aligned}
\]
the supersymmetry transformations along $Q^i_\alpha$ and $\bar Q_{\dot\alpha i}$ with commuting parameters, so that $\xi.A_m$ and $\xi.A_{kl}$ are Grassmann odd. For $\eta\in\Sigma^*$, \eqref{eq:c2-S1S} gives the special superconformal transformations: the same formulas with $(\xi_\alpha{}^i,\bar\xi_{\dot\alpha i})$ replaced by $(\epsilon_\alpha{}^i(x),\bar\epsilon_{\dot\alpha i}(x))$, and in $\eta.\lambda$ the scalar terms $\rho(A|_{W_{\mathbb R}^\perp})\eta$,
\[
\eta.\lambda_\alpha{}^k\supset-A^{kj}\varepsilon_{\alpha\beta}\eta^\beta{}_j,\qquad \eta.\bar\lambda_{\dot\beta k}\supset-A_{kj}\varepsilon_{\dot\beta\dot\alpha}\bar\eta^{\dot\alpha j}.
\]
In Example \ref{ex:c2-dil}, $\Delta.A_M=2x^m\partial_mA_M+2A_M$, $\Delta.\lambda=2x^m\partial_m\lambda+3\lambda$, $\Delta.c=2x^m\partial_mc$. In Proposition \ref{prop:c2-coeff} the spin coefficient is that of $-\tfrac18X_{mn}\sigma^m\bar\sigma^n$, the two Weyl coefficients are those of $\tfrac38\operatorname{div}v$ on $\lambda$ and $\tfrac14\operatorname{div}v$ on $A_I$, and the coefficient $-1$ is that of the scalar terms above.
\end{remark}

% ANCHOR: after_label def:c2-S2
\begin{remark}[Physics notation: the quadratic-antifield term]
For $\epsilon_1,\epsilon_2$ each a constant $\xi$ or $\rho(x)\eta$,
\[
\begin{aligned}
S^{(2)}(\epsilon_1,\epsilon_2)=\int d^4x\Big\{&\tfrac14\Big[-2i\,\Gamma^m(\epsilon_1,\epsilon_2)\big(\lambda^{*\alpha}{}_i,\sigma^m_{\alpha\dot\alpha}\bar\lambda^{*\dot\alpha i}\big)+\tfrac14\Gamma_{kl}(\epsilon_1,\epsilon_2)\Gamma^{kl}(\lambda^*,\lambda^*)\Big]-\tfrac12\big((\epsilon_1,\lambda^*),(\epsilon_2,\lambda^*)\big)\\
&-\big(\Gamma^m(\epsilon_1,\epsilon_2)A_m+\tfrac14\Gamma_{ij}(\epsilon_1,\epsilon_2)A^{ij},\,c^*\big)\Big\},
\end{aligned}
\]
with $\Gamma^m$, $\Gamma_{kl}$ the components of the $\Gamma$-pairing on $\Sigma$ evaluated on $(\epsilon_\alpha{}^i(x),\bar\epsilon_{\dot\alpha i}(x))$, $(\epsilon,\lambda^*)=\epsilon_\alpha{}^i\lambda^{*\alpha}{}_i-\bar\epsilon_{\dot\alpha i}\bar\lambda^{*\dot\alpha i}$ and
\[
\Gamma^{kl}(\lambda^*,\lambda^*)=-2\varepsilon^{klpj}\big(\varepsilon_{\alpha\beta}\lambda^{*\beta}{}_j,\lambda^{*\alpha}{}_p\big)+2\big(\varepsilon_{\dot\beta\dot\alpha}\bar\lambda^{*\dot\alpha l},\bar\lambda^{*\dot\beta k}\big)-2\big(\varepsilon_{\dot\beta\dot\alpha}\bar\lambda^{*\dot\alpha k},\bar\lambda^{*\dot\beta l}\big).
\]
In $\{S,S^{(2)}\}$ the $\lambda^*\lambda^*$ terms produce the gaugino equation of motion, and the $c^*$ term the gauge transformation with the field-dependent parameter $\Gamma^mA_m+\tfrac14\Gamma_{ij}A^{ij}=\Gamma^mA_m+\Gamma^IA_I$.
\end{remark}

% ANCHOR: end_of_section 2.4
\begin{remark}[Physics notation: gauge parameters and the end of the tower]
The gauge parameter $\iota_{\Gamma(Q_1,Q_2)}A=\Gamma^mA_m+\Gamma^IA_I$ is:
\begin{itemize}
\item $Q_1,Q_2\in\Sigma$: constant, with $\Gamma^m=i\big(\bar\xi_{2\dot\alpha i}\bar\sigma^{m\dot\alpha\alpha}\xi_{1\alpha}{}^i+\bar\xi_{1\dot\alpha i}\bar\sigma^{m\dot\alpha\alpha}\xi_{2\alpha}{}^i\big)$ and $\Gamma_{kl}(\xi_1,\xi_2)$;
\item $Q_1\in\Sigma$, $Q_2\in\Sigma^*$: the same components evaluated on $\xi_1$ and on $\rho(x)\eta_2$, linear in $x$;
\item $Q_1,Q_2\in\Sigma^*$: $-\big(x^2v_1^m-2(v_1\cdot x)x^m\big)A_m-\tfrac14x^2\,\Gamma^{kl}(\eta_1,\eta_2)A_{kl}$, with $v_1^m=-i\big(\eta_2{}^\alpha{}_i\sigma^m_{\alpha\dot\alpha}\bar\eta_1^{\dot\alpha i}+\eta_1{}^\alpha{}_i\sigma^m_{\alpha\dot\alpha}\bar\eta_2^{\dot\alpha i}\big)$.
\end{itemize}
$S^{(2)}$ is built from $A_m$, $A_I$, $\lambda^{*\alpha}{}_i$, $\bar\lambda^{*\dot\alpha i}$ and $c^*$, which pair with $A^*_M$, $\lambda_\alpha{}^i$, $\bar\lambda_{\dot\alpha i}$ and $c$, so $\{S^{(2)},S^{(2)}\}=0$; with $S^{(k\ge3)}=0$ the equivariant action is at most quadratic in the constant ghosts.
\end{remark}

% ANCHOR: end_of_section 2.5
\begin{remark}[Physics notation: closure and the ten-dimensional input]
The subalgebra of Axiom \ref{ax:c2-esw}, generated by $P_m$, the rotations, the R-symmetry, $Q^i_\alpha$ and $\bar Q_{\dot\alpha i}$, is the $\mathcal N=4$ super-Poincaré algebra with its R-symmetry, obtained by reduction from ten-dimensional super Yang--Mills theory with gaugino in the $\mathbf{16}$ of $\mathrm{Spin}(10)$. In the physical reading, for abelian free fields with $\xi.A=-\Gamma(\xi,\lambda)$, $\xi.\lambda=-\tfrac12\rho(F_A)\xi$ and $v=\Gamma(\xi_a,\xi_b)$,
\[
\{\xi_a.,\xi_b.\}A_m=-v^nF_{nm}+v^I\partial_mA_I=-v^n\partial_nA_m+\partial_m\big(v^nA_n+v^IA_I\big),\qquad \{\xi_a.,\xi_b.\}A_{ij}=-v^n\partial_nA_{ij},
\]
the translation $-v^n\partial_n$ together with the gauge transformation with parameter $\iota_vA=v^nA_n+v^IA_I$ of the $c^*$ term of $S^{(2)}$. In Lemma \ref{lem:c2-colour} the tree tensors are, for $\mathfrak{su}(N)$ and up to normalisation, products of structure constants $f^{abc}$ contracted with $\delta^{ab}$ along a tree, and $\delta^{ab}$ for $k=2$.
\end{remark}
